% ======================================================================
% Complete Proofs
% Integrated proof appendix for the TMLR submission.
% Restates each main-text result before its proof; internal lemmas are
% numbered internally and do not collide with main-text numbering.
% Main-text proof sketches are removed; this appendix supplies the complete arguments.
% ======================================================================

\section{Complete Proofs}
\label{app:proofs}

This appendix contains complete proofs of Proposition~\ref{prop:kl},
Theorems~\ref{thm:exact}, \ref{thm:hybrid}, \ref{thm:coord}, \ref{thm:fs},
\ref{thm:cmc}, \ref{thm:fixedsuite}, and~\ref{thm:population}, together with
Propositions~\ref{prop:persistent-token}, \ref{prop:gap}, \ref{prop:hazard},
\ref{prop:cmc-onestep-opt}, and~\ref{prop:lifted}, the margin-necessity lemma, and
Corollaries~\ref{cor:fs-horizon}, \ref{cor:functional}, \ref{cor:det-limit},
\ref{cor:fixedsuite}, and~\ref{cor:population}. Each main-text statement is
restated before its proof. All state spaces are finite, so no measurability issues arise and every
optimal transport problem used below attains its minimum.

\subsection{Notation and standing assumptions}
\label{app:proofs:setup}

Throughout, $\Vs$ is a finite vocabulary, $\Mtok\notin\Vs$, $\Xs=(\Vs\cup\{\Mtok\})^{N}$,
$\mathcal{M}(s)=\{i:s[i]=\Mtok\}$, and reverse time runs $t=T,\dots,0$. The metric is
$d(s,s')=\sum_{i=1}^{N}\lambda_{i}\Ind{s[i]\neq s'[i]}$ with $\lambda_{i}>0$,
$\sum_{i}\lambda_{i}=1$, so $\operatorname{diam}(d)=1$. We isolate the assumptions used by each
result:

\begin{itemize}[leftmargin=2.2em,itemsep=1pt]
\item[\textbf{(A1)}] $\Vs$, $N$, and $T$ are finite; $P_t,Q_t$ are Markov kernels on $\Xs$.
\item[\textbf{(A2)}] Common initial law: $X_T^\star\sim\mu_T$ and $X_T^q\sim\mu_T$; for
pathwise statements the coupled initial states are equal almost surely.
\item[\textbf{(A3)}] \emph{Monotone copy-only support}: for $K_t\in\{P_t,Q_t\}$, if
$K_t(s,s')>0$ and $s[i]\in\Vs$, then $s'[i]=s[i]$. Thus an unmasked token is never rewritten or
remasked. A coordinate that is masked may remain masked or become a vocabulary token.
\item[\textbf{(A4)}] \emph{Deterministic action special case}: when invoked, a map $F_t^M$ uses a
commit set $C_t^M(s)\subseteq\mathcal M(s)$, writes tokens in $\Vs$ on that set, and copies all
other coordinates.
\item[\textbf{(A5)}] \emph{Common deterministic tie-break}: when an argmax or top-$k$ statement is
invoked, both models use the same fixed strict ordering to resolve ties.
\end{itemize}

\begin{definition}[Synchronous representation]
\label{def:app:sync}
A family of maps $F_t^M:\Xs\times[0,1]\to\Xs$, $M\in\{\star,q\}$, together with independent
uniforms $U_T,\dots,U_1$ independent of $X_T$, is a synchronous representation if
$F_t^\star(s,U_t)$ has law $P_t(s,\cdot)$ and $F_t^q(s,U_t)$ has law $Q_t(s,\cdot)$. Every kernel
on a finite space admits such a representation.

For $j\in\{0,\dots,T\}$, define recursively $Y_T^{(j)}=X_T$ and, for $u=T,\dots,1$,
\[
Y_{u-1}^{(j)}=
\begin{cases}
F_u^\star(Y_u^{(j)},U_u), & u>j,\\
F_u^q(Y_u^{(j)},U_u), & u\le j.
\end{cases}
\]
The canonical pathwise hybrid output is $Z^{(j)}:=Y_0^{(j)}$. Thus $Z^{(0)}=X_0^\star$ and
$Z^{(T)}=X_0^q$ under the corresponding common-random-number coupling.
\end{definition}

\subsection{Wasserstein preliminaries}
\label{app:proofs:wass}

\begin{lemma}[Kantorovich distance on a finite space]
\label{lem:G:metric}
$\Wd(\mu,\nu)=\min_{\gamma\in\Pi(\mu,\nu)}\Ex_{\gamma}[d(X,Y)]$ is attained, and $\Wd$ is a metric
on the probability simplex over $\Xs$.
\end{lemma}

\begin{proof}
$\Pi(\mu,\nu)$ is a nonempty compact subset of a finite-dimensional simplex and
$\gamma\mapsto\Ex_{\gamma}[d]$ is linear, hence the minimum is attained. Symmetry and
$\Wd(\mu,\nu)=0\iff\mu=\nu$ are immediate from $d$ being a metric. For the triangle inequality,
let $\gamma_{12}\in\Pi(\mu_{1},\mu_{2})$ and $\gamma_{23}\in\Pi(\mu_{2},\mu_{3})$ be optimal.
Define the glued law on $\Xs^{3}$ by
$\gamma(x,y,z)=\gamma_{12}(x,y)\,\gamma_{23}(y,z)/\mu_{2}(y)$ when $\mu_{2}(y)>0$ and $0$
otherwise; its $(1,2)$ and $(2,3)$ marginals are $\gamma_{12}$ and $\gamma_{23}$, so its $(1,3)$
marginal lies in $\Pi(\mu_{1},\mu_{3})$, and
$\Ex_{\gamma}[d(X,Z)]\le\Ex_{\gamma}[d(X,Y)]+\Ex_{\gamma}[d(Y,Z)]
=\Wd(\mu_{1},\mu_{2})+\Wd(\mu_{2},\mu_{3})$.
\end{proof}

\begin{lemma}[Convexity under a common input law]
\label{lem:G:convex}
For any probability law $\mu$ and kernels $K_{1},K_{2}$,
$\Wd(\mu K_{1},\mu K_{2})\le\sum_{s}\mu(s)\,\Wd\big(K_{1}(s,\cdot),K_{2}(s,\cdot)\big)$.
\end{lemma}

\begin{proof}
For each $s$ pick an optimal $\gamma_{s}\in\Pi(K_{1}(s,\cdot),K_{2}(s,\cdot))$
(Lemma~\ref{lem:G:metric}) and set $\gamma=\sum_{s}\mu(s)\gamma_{s}$. Its marginals are
$\mu K_{1}$ and $\mu K_{2}$, and its cost is the right-hand side; conclude by minimality.
\end{proof}

\subsection{Proof of Proposition~\ref{prop:kl} (path-KL vacuity)}
\label{app:proofs:kl}

\paragraph{Proposition~\ref{prop:kl} (restated).}
\emph{If there exists a teacher-reachable state $s$ at some step $t$ (i.e.\ $\mu^{\star}_{t}(s)>0$)
with $P_{t}(s,\cdot)\not\ll Q_{t}(s,\cdot)$, then
$\mathrm{KL}\big(\mathbb{P}^{\star}\,\big\|\,\mathbb{P}^{q}\big)=+\infty$ for the trajectory laws
on $\Xs^{T+1}$. In particular this holds when deterministic decoders choose different next states
at $s$; and, under a deterministic commit-set selector with stochastic token sampling, whenever
$C^{\star}_{t}(s)\neq C^{q}_{t}(s)$.}

\begin{proof}
By hypothesis there is $s'$ with $P_{t}(s,s')>0=Q_{t}(s,s')$. Consider the trajectory event
$E=\{(x_{T},\dots,x_{0}):x_{t}=s,\ x_{t-1}=s'\}$. By the Markov property,
\[
\mathbb{P}^{\star}(E)=\mu^{\star}_{t}(s)\,P_{t}(s,s')>0,
\qquad
\mathbb{P}^{q}(E)=\mu^{q}_{t}(s)\,Q_{t}(s,s')=0 .
\]
Hence $\mathbb{P}^{\star}\not\ll\mathbb{P}^{q}$, and by the definition of KL divergence on a finite
space (with the convention $a\log(a/0)=+\infty$ for $a>0$),
$\mathrm{KL}(\mathbb{P}^{\star}\|\mathbb{P}^{q})=+\infty$.

For the deterministic sub-case, $P_{t}(s,\cdot)=\delta_{F^{\star}_{t}(s)}$ and
$Q_{t}(s,\cdot)=\delta_{F^{q}_{t}(s)}$ with $F^{\star}_{t}(s)\neq F^{q}_{t}(s)$; take
$s'=F^{\star}_{t}(s)$. For the stochastic-token sub-case, assume each model first chooses a deterministic commit set
$C_t^M(s)\subseteq\mathcal M(s)$, writes only vocabulary tokens on that set, and copies the
remaining coordinates. Every state in the support of $P_t(s,\cdot)$ then has mask set exactly
$\mathcal M(s)\setminus C_t^\star(s)$, while every state in the support of $Q_t(s,\cdot)$ has mask
set $\mathcal M(s)\setminus C_t^q(s)$. Distinct commit sets therefore give disjoint supports.

Finally, the marginalization map $(x_{T},\dots,x_{0})\mapsto x_{0}$ is a deterministic
pushforward, and KL divergence is nonincreasing under pushforwards (data-processing inequality),
so $\mathrm{KL}(\mu^{\star}_{0}\|\mu^{q}_{0})\le\mathrm{KL}(\mathbb{P}^{\star}\|\mathbb{P}^{q})$.
An infinite right-hand side therefore carries no information about the output KL, which is the
one-directional reading stated in the main text.
\end{proof}

\subsection{Proof of Theorem~\ref{thm:exact} (exact one-step decomposition)}
\label{app:proofs:action}

\paragraph{Theorem~\ref{thm:exact} (restated).}
\emph{For monotone copy-only deterministic decoders evaluated at a shared state $s$, with commit
sets $C^{\star}_{t},C^{q}_{t}\subseteq\mathcal{M}(s)$ and committed tokens
$a^{\star}_{t,i},a^{q}_{t,i}\in\Vs$,}
\[
\epsilon_{t}(s)\;=\;d\big(F^{\star}_{t}(s),F^{q}_{t}(s)\big)
\;=\;\sum_{i\in C^{\star}_{t}\triangle C^{q}_{t}}\lambda_{i}
\;+\;\sum_{i\in C^{\star}_{t}\cap C^{q}_{t}}\lambda_{i}\,\Ind{a^{\star}_{t,i}\neq a^{q}_{t,i}} .
\]

\begin{proof}
Write $s^{\star}=F^{\star}_{t}(s)$, $s^{q}=F^{q}_{t}(s)$ and partition $\{1,\dots,N\}$ into four
classes.

\emph{(i) $i\notin\mathcal{M}(s)$.} By \textbf{(A4)} the position is in neither commit set and both decoders copy it, so $s^{\star}[i]=s^{q}[i]=s[i]$ and the contribution is $0$.

\emph{(ii) $i\in\mathcal{M}(s)\setminus(C^{\star}_{t}\cup C^{q}_{t})$.} Both copy the mask:
$s^{\star}[i]=s^{q}[i]=\Mtok$; contribution $0$.

\emph{(iii) $i\in C^{\star}_{t}\triangle C^{q}_{t}$.} Exactly one decoder commits. The committing
side writes a token in $\Vs$ by \textbf{(A4)}; the other side copies $\Mtok$. Since
$\Mtok\notin\Vs$, $s^{\star}[i]\neq s^{q}[i]$ \emph{whatever} token was written; contribution
$\lambda_{i}$.

\emph{(iv) $i\in C^{\star}_{t}\cap C^{q}_{t}$.} Both commit; $s^{\star}[i]=a^{\star}_{t,i}$ and
$s^{q}[i]=a^{q}_{t,i}$, so the contribution is $\lambda_{i}\Ind{a^{\star}_{t,i}\neq a^{q}_{t,i}}$.

Summing the four classes against $\lambda$ gives the identity.
\end{proof}

\subsection{Proof of Proposition~\ref{prop:persistent-token} (persistent token lower bound)}
\label{app:proofs:persistence}

\paragraph{Proposition~\ref{prop:persistent-token} (restated).}
\emph{For deterministic actions at a shared state, let
$I=\{i\in C_t^\star\cap C_t^q:a_{t,i}^\star\ne a_{t,i}^q\}$. Under a common monotone copy-only
suffix, $\eta_t(s)\ge\sum_{i\in I}\lambda_i$.}

\begin{proof}
Let $s^\star=F_t^\star(s)$ and $s^q=F_t^q(s)$. For every $i\in I$,
$s^\star[i]=a_{t,i}^\star\ne a_{t,i}^q=s^q[i]$, and both values are vocabulary tokens. By
Assumption~\textbf{(A3)}, every terminal state in the support of
$\delta_{s^\star}Q_{t-1:1}$ retains $a_{t,i}^\star$ at coordinate $i$, while every terminal state
in the support of $\delta_{s^q}Q_{t-1:1}$ retains $a_{t,i}^q$. Hence every coupling of these two
terminal laws has cost at least $\sum_{i\in I}\lambda_i$. Taking the infimum over couplings gives
the claimed lower bound on $\eta_t(s)$.
\end{proof}

\subsection{Proof of Lemma~\ref{lem:margin-main} (margin necessity)}
\label{app:proofs:margins}

We now prove the two inclusions stated in Lemma~\ref{lem:margin-main}. Assumption~\textbf{(A5)} handles ties: fix a strict total order $\prec$ on indices, and let both decoders select
the $\prec$-least maximizer.

\begin{lemma}[Token-margin necessity]
\label{lem:G:tokmargin}
Fix a position $i$ and logit vectors $z^{\star},z^{q}\in\mathbb{R}^{|\Vs|}$ with
$\|z^{q}-z^{\star}\|_{\infty}\le R$. Let $a$ and $b$ be the tie-broken argmaxes of $z^{\star}$ and
$z^{q}$ respectively, and let
$M^{\mathrm{tok}}=z^{\star}_{(1)}-z^{\star}_{(2)}$ be the top-1/top-2 gap of $z^{\star}$. If
$b\neq a$ then $M^{\mathrm{tok}}\le2R$.
\end{lemma}

\begin{proof}
Since $b$ is selected under $z^{q}$ over $a$, either $z^{q}_{b}>z^{q}_{a}$, or
$z^{q}_{b}=z^{q}_{a}$ with $b\prec a$. In the first case,
\[
z^{\star}_{a}-z^{\star}_{b}
=(z^{\star}_{a}-z^{q}_{a})+(z^{q}_{a}-z^{q}_{b})+(z^{q}_{b}-z^{\star}_{b})
\le R+0+R=2R .
\]
In the second case, $a$ being selected under $z^{\star}$ over $b$ while $b\prec a$ forces
$z^{\star}_{a}>z^{\star}_{b}$, and
$z^{\star}_{a}-z^{\star}_{b}
=(z^{\star}_{a}-z^{q}_{a})+(z^{q}_{b}-z^{\star}_{b})\le2R$
using $z^{q}_{a}=z^{q}_{b}$. In both cases, since $b\neq a$ we have
$z^{\star}_{b}\le z^{\star}_{(2)}$, whence
$M^{\mathrm{tok}}=z^{\star}_{(1)}-z^{\star}_{(2)}\le z^{\star}_{a}-z^{\star}_{b}\le2R$.
\end{proof}

\begin{lemma}[Policy-margin necessity]
\label{lem:G:polmargin}
Let $c^{\star},c^{q}\in\mathbb{R}^{\mathcal{M}(s)}$ be confidence scores with
$\max_{i}|c^{q}_{i}-c^{\star}_{i}|\le R^{\mathrm{pol}}$, and let $C^{\star},C^{q}$ be the
tie-broken top-$k$ index sets under $c^{\star}$ and $c^{q}$ respectively, with
$k<|\mathcal{M}(s)|$. Let
$M^{\mathrm{pol}}=c^{\star}_{(k)}-c^{\star}_{(k+1)}$ (order statistics of $c^{\star}$). If
$C^{q}\neq C^{\star}$ then $M^{\mathrm{pol}}\le2R^{\mathrm{pol}}$. (If $|\mathcal{M}(s)|\le k$,
both decoders commit all masked positions and no policy split can occur; by convention
$M^{\mathrm{pol}}=+\infty$.)
\end{lemma}

\begin{proof}
Since $|C^{\star}|=|C^{q}|=k$ and $C^{q}\neq C^{\star}$, there exist
$i\in C^{\star}\setminus C^{q}$ and $j\in C^{q}\setminus C^{\star}$. Membership under $c^{\star}$
gives $c^{\star}_{i}\ge c^{\star}_{(k)}$ and $c^{\star}_{j}\le c^{\star}_{(k+1)}$. Selection of
$j$ over $i$ under $c^{q}$ gives $c^{q}_{j}>c^{q}_{i}$, or $c^{q}_{j}=c^{q}_{i}$ with
$j\prec i$; in the latter case, $i$'s selection over $j$ under $c^{\star}$ forces
$c^{\star}_{i}>c^{\star}_{j}$, and in either case, exactly as in
Lemma~\ref{lem:G:tokmargin},
$c^{\star}_{i}-c^{\star}_{j}\le2R^{\mathrm{pol}}$. Therefore
$M^{\mathrm{pol}}\le c^{\star}_{i}-c^{\star}_{j}\le2R^{\mathrm{pol}}$.
\end{proof}

Both inclusions are necessary conditions only; no claim of sufficiency is made or used anywhere in
the paper.

\subsection{Proof of Theorem~\ref{thm:hybrid} (hybrid propagation)}
\label{app:proofs:hybrid}

\paragraph{Theorem~\ref{thm:hybrid} (restated).}
\emph{Under \textbf{(A1)}--\textbf{(A2)}, with
$\eta_{t}(s)=\Wd\big(P_{t}(s,\cdot)Q_{t-1:1},\,Q_{t}(s,\cdot)Q_{t-1:1}\big)$,}
\[
\Wd(\mu^{\star}_{0},\mu^{q}_{0})\ \le\ \sum_{t=1}^{T}\Ex_{S_{t}\sim\mu^{\star}_{t}}\big[\eta_{t}(S_{t})\big].
\]
\emph{Moreover, for the CRN coupling of Definition~\ref{def:app:sync}, the pathwise analogue
$\Ex[d(X^{\star}_{0},X^{q}_{0})]\le\sum_{t}\Ex[d(Z^{(t-1)},Z^{(t)})]$ holds.}

\begin{proof}
Define the hybrid \emph{laws} $H_{j}=\mu_{T}P_{T}\cdots P_{j+1}Q_{j}\cdots Q_{1}$ for
$j=0,\dots,T$ (empty products are identities), so $H_{0}=\mu^{\star}_{0}$ and
$H_{T}=\mu^{q}_{0}$. By Lemma~\ref{lem:G:metric} and the triangle inequality along the chain
$H_{0},H_{1},\dots,H_{T}$,
\[
\Wd(\mu^{\star}_{0},\mu^{q}_{0})\ \le\ \sum_{j=1}^{T}\Wd(H_{j-1},H_{j}).
\]
Fix $j$. Both $H_{j-1}$ and $H_{j}$ push the same prefix law
$\mu_{T}P_{T}\cdots P_{j+1}=\mu^{\star}_{j}$ through a suffix kernel: writing
$K_{1}=P_{j}Q_{j-1:1}$ and $K_{2}=Q_{j}Q_{j-1:1}$,
\[
H_{j-1}=\mu^{\star}_{j}K_{1},\qquad H_{j}=\mu^{\star}_{j}K_{2}.
\]
Lemma~\ref{lem:G:convex} then gives
$\Wd(H_{j-1},H_{j})\le\Ex_{S_{j}\sim\mu^{\star}_{j}}
\big[\Wd\big(K_{1}(S_{j},\cdot),K_{2}(S_{j},\cdot)\big)\big]
=\Ex_{\mu^{\star}_{j}}[\eta_{j}]$. Summing over $j$ proves the distributional bound.

For the pathwise analogue, note that $(Z^{(j)})_{j=0}^{T}$ of Definition~\ref{def:app:sync} are
$T+1$ random states on one probability space with $Z^{(0)}=X^{\star}_{0}$ and
$Z^{(T)}=X^{q}_{0}$. The pointwise triangle inequality of the metric $d$ gives, for every sample
path,
$d(Z^{(0)},Z^{(T)})\le\sum_{t=1}^{T}d(Z^{(t-1)},Z^{(t)})$; take expectations. (Each pair
$(Z^{(t-1)},Z^{(t)})$ is, marginally, a coupling of $H_{t-1}$ and $H_{t}$, so the pathwise sum
dominates the distributional one term by term; the two versions are consistent.)
\end{proof}

\begin{remark}[The occupancy--suffix asymmetry is forced]
\label{rem:G:asym}
In the display above, the step-$j$ input law is the \emph{teacher} occupancy $\mu^{\star}_{j}$
while the propagation kernel $Q_{j-1:1}$ is the \emph{quantized} suffix, because consecutive
hybrids share the FP prefix and the Q suffix by construction. Reversing the order of the
telescoping (Q prefix, P suffix) yields the mirror pairing
$\Ex_{\mu^{q}_{j}}[\Wd(Q_{j}P_{j-1:1},P_{j}P_{j-1:1})]$; no ordering yields teacher occupancy
paired with teacher propagation. This is the precise content of the claim that the two chains
enter asymmetrically.
\end{remark}

\subsection{Proof of Theorem~\ref{thm:coord} (coordinate-union hierarchy)}
\label{app:proofs:coord}

Throughout this subsection fix the CRN synchronous representation and the canonical hybrids
$Z^{(0)},\dots,Z^{(T)}$ of Definition~\ref{def:app:sync}; set
$A_{t,i}=\Ind{Z^{(t-1)}[i]\neq Z^{(t)}[i]}$, $N_{i}=\sum_{t=1}^{T}A_{t,i}$,
$q_{t,i}=\Ex[A_{t,i}]$, and
\begin{align*}
U_{\mathrm{union}}(Z)&=\sum_i\lambda_i\Ind{N_i\ge1},
& B_{\mathrm{union}}&=\Ex[U_{\mathrm{union}}(Z)],\\
B_{\mathrm{coord}}&=\sum_i\lambda_i\min\{1,\Ex[N_i]\},
& B_{\mathrm{tel}}&=\sum_t\Ex\big[d(Z^{(t-1)},Z^{(t)})\big].
\end{align*}

\paragraph{Theorem~\ref{thm:coord} (restated).}
$\displaystyle
\Wd(\mu^{\star}_{0},\mu^{q}_{0})\ \le\ \Dcrn\ \le\ B_{\mathrm{union}}\ \le\ B_{\mathrm{coord}}
\ \le\ \min\{1,\,B_{\mathrm{tel}}\}.$

\begin{proof}
\emph{First inequality.} $(Z^{(0)},Z^{(T)})=(X^{\star}_{0},X^{q}_{0})$ is a coupling of
$(\mu^{\star}_{0},\mu^{q}_{0})$: $Z^{(0)}$ applies only full-precision maps to $X_{T}$ with the
uniforms, hence has law $\mu^{\star}_{0}$, and symmetrically for $Z^{(T)}$. Minimality of $\Wd$
over couplings gives $\Wd\le\Dcrn$.

\emph{Second inequality.} Fix a sample path and a coordinate $i$. If $A_{t,i}=0$ for every $t$,
then $Z^{(0)}[i]=Z^{(1)}[i]=\dots=Z^{(T)}[i]$, so
$\Ind{Z^{(0)}[i]\neq Z^{(T)}[i]}\le\Ind{N_{i}\ge1}$ pointwise. Multiplying by $\lambda_{i}$ and
summing,
$d(X^{\star}_{0},X^{q}_{0})=d(Z^{(0)},Z^{(T)})\le U(Z)$ pointwise; take expectations.

\emph{Third inequality.} For each $i$, $\Pr[N_{i}\ge1]\le1$ trivially and
$\Pr[N_{i}\ge1]\le\Ex[N_{i}]$ by Markov's inequality applied to the nonnegative integer variable
$N_{i}$; hence $\Pr[N_{i}\ge1]\le\min\{1,\Ex[N_{i}]\}$. Multiply by $\lambda_{i}$ and sum.

\emph{Fourth inequality.} Termwise $\min\{1,\Ex[N_{i}]\}\le1$ gives
$B_{\mathrm{coord}}\le\sum_{i}\lambda_{i}=1$, and termwise
$\min\{1,\Ex[N_{i}]\}\le\Ex[N_{i}]$ gives
\[
B_{\mathrm{coord}}\ \le\ \sum_{i}\lambda_{i}\,\Ex[N_{i}]
=\sum_{i}\lambda_{i}\sum_{t}q_{t,i}
=\sum_{t}\Ex\Big[\sum_{i}\lambda_{i}A_{t,i}\Big]
=\sum_{t}\Ex\big[d(Z^{(t-1)},Z^{(t)})\big]=B_{\mathrm{tel}},
\]
where the middle equality is Tonelli for finite sums and the next uses the definition of $d$.
\end{proof}

\label{app:proofs:gap}
\begin{proposition}[Exact gap decomposition]
\label{prop:G:gap}
\[
B_{\mathrm{coord}}-\Dcrn
=\sum_{i}\lambda_{i}\Big[
\underbrace{\min\{1,\Ex[N_{i}]\}-\Pr[N_{i}\ge1]}_{\text{multiplicity gap }\ge0}
+\underbrace{\Pr[N_{i}\ge1]-\Pr\big[Z^{(0)}[i]\neq Z^{(T)}[i]\big]}_{\text{cancellation gap }\ge0}
\Big].
\]
Moreover, when $\Ex[N_{i}]\le1$ the multiplicity gap equals $\Ex[(N_{i}-1)^{+}]$, so it vanishes
iff $N_{i}\in\{0,1\}$ a.s.; and the cancellation gap vanishes iff, for every coordinate
$i$ with $\lambda_i>0$, $\Pr[N_i\ge1,\,Z^{(0)}[i]=Z^{(T)}[i]]=0$ (equivalently, a changed
coordinate ends different almost surely).
\end{proposition}

\begin{proof}
By the terminal identity $\Dcrn=\sum_{i}\lambda_{i}\Pr[Z^{(0)}[i]\neq Z^{(T)}[i]]$ (definition of
$d$ plus linearity), the display telescopes. Nonnegativity of the two brackets is the third and
second inequalities of Theorem~\ref{thm:coord} respectively. If $\Ex[N_{i}]\le1$ then
$\min\{1,\Ex[N_{i}]\}-\Pr[N_{i}\ge1]=\Ex[N_{i}]-\Pr[N_{i}\ge1]
=\sum_{k\ge1}(k-1)\Pr[N_{i}=k]=\Ex[(N_{i}-1)^{+}]$.
\end{proof}

\subsection{Proof of Theorem~\ref{thm:fs} (first-divergence certificate)}
\label{app:proofs:fs}

\begin{definition}[First divergence]
\label{def:G:tau}
Let $\Gamma$ be any coupling of the two chains with $X^{\star}_{T}=X^{q}_{T}$. If the chains ever
split, let $\tau$ be the unique $t\in\{1,\dots,T\}$ such that
$X^{\star}_{u}=X^{q}_{u}$ for all $u\ge t$ and $X^{\star}_{t-1}\neq X^{q}_{t-1}$; otherwise set
$\tau=\dagger$. Uniqueness follows because any earlier chronological split would violate the
agreement condition for all $u\ge t$. On $\{\tau=t\}$ the chains agree at the input of step $t$;
write $X_{\tau}:=X^{\star}_{\tau}=X^{q}_{\tau}$. Later re-agreement is permitted and does not alter
$\tau$.
\end{definition}

\begin{lemma}[Frozen prefix]
\label{lem:G:frozen}
Assume \textbf{(A3)} for each chain separately. If $X^{\star}_{t}=X^{q}_{t}=s$,
then for every $u\le t$ and every $i\notin\mathcal{M}(s)$,
$X^{\star}_{u}[i]=X^{q}_{u}[i]=s[i]$.
\end{lemma}

\begin{proof}
Fix a chain $M\in\{\star,q\}$ and argue by downward induction on $u$ from $u=t$. Base case:
$X^{M}_{t}[i]=s[i]$. Inductive step: if $X^{M}_{u}[i]=s[i]\in\Vs$ (position $i$ is not masked in $s$), Assumption~\textbf{(A3)} gives
$X^{M}_{u-1}[i]=X^{M}_{u}[i]=s[i]$. The two chains agree at $i$ for all $u\le t$ because each separately equals $s[i]$.
\end{proof}

\paragraph{Theorem~\ref{thm:fs} (restated).}
\emph{For any coupling $\Gamma$ of monotone copy-only chains from a common initial state, with
$R(s)=\sum_{i\in\mathcal{M}(s)}\lambda_{i}$,}
\[
\Wd(\mu^{\star}_{0},\mu^{q}_{0})\ \le\ D_{\Gamma}\ \le\
\Bfs^{\Gamma}:=\Ex_{\Gamma}\big[R(X_{\tau})\Ind{\tau\neq\dagger}\big]\ \le\
\Pr_{\Gamma}[\tau\neq\dagger]\ \le\ 1 .
\]

\begin{proof}
The first inequality is minimality of $\Wd$ over couplings. For the second, condition on the value
of $\tau$. On $\{\tau=\dagger\}$ the chains coincide at every time, in particular
$d(X^{\star}_{0},X^{q}_{0})=0$. On $\{\tau=t\}$, apply Lemma~\ref{lem:G:frozen} with
$s=X_{\tau}$: for every $i\notin\mathcal{M}(X_{\tau})$ we have
$X^{\star}_{0}[i]=X^{q}_{0}[i]$, hence
\[
d(X^{\star}_{0},X^{q}_{0})
=\sum_{i}\lambda_{i}\Ind{X^{\star}_{0}[i]\neq X^{q}_{0}[i]}
=\sum_{i\in\mathcal{M}(X_{\tau})}\lambda_{i}\Ind{X^{\star}_{0}[i]\neq X^{q}_{0}[i]}
\le R(X_{\tau}).
\]
Taking expectations over the disjoint events $\{\tau=\dagger\},\{\tau=T\},\dots,\{\tau=1\}$
yields $D_{\Gamma}\le\Bfs^{\Gamma}$. The last two inequalities follow from $R\le1$ and
probabilities being at most $1$.
\end{proof}

\subsection{Proof of Corollary~\ref{cor:fs-horizon} (finite-horizon first divergence)}
\label{app:proofs:fs-horizon}

\paragraph{Corollary~\ref{cor:fs-horizon} (restated).}
\emph{For $r\in\{0,\dots,T-1\}$, stop the coupled chains at time $r$ and let
$\tau^{(r)}\in\{r+1,\dots,T\}\cup\{\dagger\}$ be their unique first divergence before that
stopping time. Then}
\[
\Wd(\mu_r^\star,\mu_r^q)
\le \Ex_\Gamma[d(X_r^\star,X_r^q)]
\le \Ex_\Gamma[R(X_{\tau^{(r)}})\Ind{\tau^{(r)}\neq\dagger}]
\le1.
\]

\begin{proof}
The first inequality is the defining minimality of Wasserstein distance over couplings of
$(\mu_r^\star,\mu_r^q)$. If $\tau^{(r)}=\dagger$, the two states agree at time $r$. On
$\{\tau^{(r)}=t\}$, Lemma~\ref{lem:G:frozen} applied at the common pre-split state
$X_t^\star=X_t^q$ shows that every coordinate already unmasked at time $t$ remains identical in
both chains down to time $r$. Hence only coordinates in $\mathcal M(X_t)$ can contribute to
$d(X_r^\star,X_r^q)$, and the same pointwise calculation as in Theorem~\ref{thm:fs} gives
$d(X_r^\star,X_r^q)\le R(X_t)$. Averaging over the disjoint first-divergence events proves the
claim. Replacing the global start time $T$ by a deterministic window start at which the coupled
states agree gives the windowed statement verbatim.
\end{proof}



\subsection{Proof of Corollary~\ref{cor:functional} (bounded-functional stability)}
\label{app:proofs:functional}

\paragraph{Corollary~\ref{cor:functional} (restated).}
\emph{Under the setting of Corollary~\ref{cor:fs-horizon}, for every fixed
$f:\Xs\to[0,1]$,
$|\Ex f(X_r^\star)-\Ex f(X_r^q)|
\le\Ex_\Gamma|f(X_r^\star)-f(X_r^q)|
\le\Pr_\Gamma(X_r^\star\ne X_r^q)
\le\Pr_\Gamma(\tau^{(r)}\ne\dagger)$.}

\begin{proof}
Because $\Gamma$ has marginals $\mu_r^\star$ and $\mu_r^q$,
\[
\Ex f(X_r^\star)-\Ex f(X_r^q)
=\Ex_\Gamma\!\left[f(X_r^\star)-f(X_r^q)\right].
\]
Taking absolute values and applying Jensen's inequality gives the first bound. Since $f$ takes
values in $[0,1]$, $|f(x)-f(y)|\le\Ind{x\ne y}$, which gives the second. Finally, if
$\tau^{(r)}=\dagger$, the chains have agreed throughout the window down to time $r$, so
$X_r^\star=X_r^q$. Hence
$\{X_r^\star\ne X_r^q\}\subseteq\{\tau^{(r)}\ne\dagger\}$, proving the last inequality.
For a correctness indicator, the two expectations are the corresponding accuracies.
\end{proof}

\label{app:proofs:hazard}
\begin{proposition}[Survival--hazard identity]
\label{prop:G:hazard}
Suppose in addition that $\Gamma$ admits a synchronous representation
(Definition~\ref{def:app:sync}) with one-step maps $(G^{\star}_{t},G^{q}_{t})$ driven by
$U_{t}$, and let $G_{t}=\{X^{\star}_{u}=X^{q}_{u}\ \forall u\ge t\}$ and
$h^{\Gamma}_{t}(s)=\Pr\big[G^{\star}_{t}(s,U_{t})\neq G^{q}_{t}(s,U_{t})\big]$. Then
\[
\Bfs^{\Gamma}=\sum_{t=1}^{T}\Ex\Big[\Ind{G_{t}}\;h^{\Gamma}_{t}(X_{t})\;R(X_{t})\Big].
\]
\end{proposition}

\begin{proof}
The events $\{\tau=t\}=G_{t}\cap\{G^{\star}_{t}(X_{t},U_{t})\neq G^{q}_{t}(X_{t},U_{t})\}$ are
pairwise disjoint with union $\{\tau\neq\dagger\}$, and $X_{\tau}=X_{t}$ on $\{\tau=t\}$, so
\[
\Bfs^{\Gamma}
=\sum_{t}\Ex\Big[\Ind{G_{t}}\,R(X_{t})\,
\Ind{G^{\star}_{t}(X_{t},U_{t})\neq G^{q}_{t}(X_{t},U_{t})}\Big].
\]
Both $\Ind{G_{t}}$ and $X_{t}$ are measurable with respect to
$\mathcal{F}_{t}:=\sigma(X_{T},U_{T},\dots,U_{t+1})$, and $U_{t}$ is independent of
$\mathcal{F}_{t}$. Conditioning each summand on $\mathcal{F}_{t}$ and integrating over $U_{t}$
replaces the last indicator by $h^{\Gamma}_{t}(X_{t})$.
\end{proof}

\begin{remark}[The certificate depends only on the pre-split coupling]
\label{rem:G:presplit}
$(\tau,X_{\tau})$ is measurable with respect to the trajectory up to the first split, hence
$\Bfs^{\Gamma}$ is unchanged by any modification of the coupling \emph{after} the split. In
contrast $D_{\Gamma}$ does depend on the continuation; the implemented continuation is part of
the coupling specification and the chain
$\Wd\le D_{\Gamma}\le\Bfs^{\Gamma}$ holds for that specification.
\end{remark}

\subsection{Proof of Theorem~\ref{thm:cmc} (coordinatewise maximal coupling)}
\label{app:proofs:cmc}

\begin{lemma}[Maximal coupling of two categorical laws]
\label{lem:G:maxcoupling}
Let $p,q$ be probability mass functions on a finite set $\mathcal{Y}$, $r=p\wedge q$,
$\alpha=\sum_{y}r(y)$, and $\delta=\TV(p,q)=1-\alpha$. Define a pair $(X,Y)$ by: with probability
$\alpha$, draw $X=Y\sim r/\alpha$; otherwise draw $X\sim(p-r)/\delta$ and $Y\sim(q-r)/\delta$
independently. Then $X\sim p$, $Y\sim q$, and $\Pr[X\neq Y]=\delta$, which is minimal among all
couplings of $(p,q)$. (If $\delta=0$ the residual branch is never taken; if $\alpha=0$ the common
branch is never taken.)
\end{lemma}

\begin{proof}
Marginals: for any $y$,
$\Pr[X=y]=\alpha\cdot\frac{r(y)}{\alpha}+\delta\cdot\frac{p(y)-r(y)}{\delta}=p(y)$, and
symmetrically for $Y$. Disagreement: on the common branch $X=Y$. On the residual branch the
supports $\{p>q\}$ of $p-r$ and $\{q>p\}$ of $q-r$ are disjoint, so $X\neq Y$ almost surely;
hence $\Pr[X\neq Y]=\delta$. Minimality: for any coupling $\gamma$ of $(p,q)$ and any $y$,
$\gamma(X=Y=y)\le\min\{p(y),q(y)\}=r(y)$, so $\Pr_{\gamma}[X=Y]\le\alpha$, i.e.\
$\Pr_{\gamma}[X\neq Y]\ge\delta$.
\end{proof}

\paragraph{Theorem~\ref{thm:cmc} (restated).}
\emph{Assume the one-step transitions factorize at every state:
$P_{t}(s,\cdot)=\bigotimes_{i}p^{\star}_{t,i}(\cdot\mid s)$ and
$Q_{t}(s,\cdot)=\bigotimes_{i}p^{q}_{t,i}(\cdot\mid s)$, and let
$\delta_{t,i}(s)=\TV(p^{\star}_{t,i},p^{q}_{t,i})$. Coupling each coordinate by
Lemma~\ref{lem:G:maxcoupling}, independently across coordinates, defines a one-step coupling with
both marginals exact and split probability
$h^{\mathrm{CMC}}_{t}(s)=1-\prod_{i}\big(1-\delta_{t,i}(s)\big)$, satisfying}
\[
\max_{i}\delta_{t,i}(s)\ \le\ \TV\big(P_{t}(s,\cdot),Q_{t}(s,\cdot)\big)\ \le\
h^{\mathrm{CMC}}_{t}(s)\ \le\ \min\Big\{1,\ \sum_{i}\delta_{t,i}(s)\Big\},
\]
\emph{and iterating these one-step couplings from a common initial state (with any fixed valid
continuation after the first split) yields a path coupling $\Gamma_{\mathrm{CMC}}$ with}
$\Wd(\mu^{\star}_{0},\mu^{q}_{0})\le D_{\mathrm{CMC}}\le\Bfs^{\mathrm{CMC}}\le1$.

\begin{proof}
\emph{Validity of the one-step coupling.} Let $(X_{i},Y_{i})_{i}$ be independent across $i$, each
distributed as in Lemma~\ref{lem:G:maxcoupling} for the pair
$(p^{\star}_{t,i},p^{q}_{t,i})$. Then $X=(X_{i})_{i}$ has law
$\bigotimes_{i}p^{\star}_{t,i}=P_{t}(s,\cdot)$ and $Y=(Y_{i})_{i}$ has law $Q_{t}(s,\cdot)$, by
independence and the per-coordinate marginal property. Both joint marginals are therefore
\emph{exact}; no approximation is involved.

\emph{Split probability.} $\{X\neq Y\}=\bigcup_{i}\{X_{i}\neq Y_{i}\}$, and by independence
$\Pr[X=Y]=\prod_{i}(1-\delta_{t,i})$, giving
$h^{\mathrm{CMC}}_{t}(s)=1-\prod_{i}(1-\delta_{t,i}(s))$.

\emph{Sandwich.} Lower bound: the coordinate projection $\pi_{i}$ is a deterministic map, and
total variation is nonincreasing under pushforwards, so
$\delta_{t,i}=\TV\big(\pi_{i}\#P_{t},\pi_{i}\#Q_{t}\big)\le\TV(P_{t},Q_{t})$ for every $i$; take
the maximum. Middle: $\TV$ equals the minimal disagreement probability over couplings
(Lemma~\ref{lem:G:maxcoupling} applied to the joint laws, whose proof of minimality used nothing
about $\mathcal{Y}$), and the product construction is one particular coupling, so
$\TV(P_{t},Q_{t})\le h^{\mathrm{CMC}}_{t}$. Upper: $h^{\mathrm{CMC}}\le1$ trivially, and by the
Weierstrass product inequality $\prod_{i}(1-\delta_{i})\ge1-\sum_{i}\delta_{i}$, so
$h^{\mathrm{CMC}}\le\sum_{i}\delta_{i}$.

\emph{Path coupling.} Construct $(X^{\star}_{t-1},X^{q}_{t-1})$ from
$(X^{\star}_{t},X^{q}_{t})$ as follows: if $X^{\star}_{t}=X^{q}_{t}=s$, use the one-step CMC
coupling at $s$; otherwise use any fixed coupling of
$P_{t}(X^{\star}_{t},\cdot)$ and $Q_{t}(X^{q}_{t},\cdot)$ (e.g.\ the event-addressed independent
residual coupling used in the implementation). By induction on $t$: conditionally on
$(X^{\star}_{t},X^{q}_{t})$, the first component has conditional law
$P_{t}(X^{\star}_{t},\cdot)$ and the second $Q_{t}(X^{q}_{t},\cdot)$ in both branches; hence the
first marginal path is the FP chain and the second the quantized chain. This is a valid coupling
$\Gamma_{\mathrm{CMC}}$ with a synchronous representation (each branch is a measurable function
of the current pair and fresh uniforms), so Theorem~\ref{thm:fs} applies verbatim and yields
$\Wd\le D_{\mathrm{CMC}}\le\Bfs^{\mathrm{CMC}}\le1$. By Remark~\ref{rem:G:presplit},
$\Bfs^{\mathrm{CMC}}$ does not depend on the choice of continuation, while $D_{\mathrm{CMC}}$ is
the realized cost of the implemented continuation.
\end{proof}

\subsection{Proof of Proposition~\ref{prop:cmc-onestep-opt} (one-step Wasserstein optimality)}
\label{app:proofs:cmc-opt}

\paragraph{Proposition~\ref{prop:cmc-onestep-opt} (restated).}
\emph{Under the factorization assumptions of Theorem~\ref{thm:cmc},}
\[
\Wd\big(P_t(s,\cdot),Q_t(s,\cdot)\big)=\sum_i\lambda_i\delta_{t,i}(s),
\]
\emph{and the independent coordinatewise maximal coupling attains the minimum.}

\begin{proof}
Let $\gamma$ be any coupling of $P_t(s,\cdot)$ and $Q_t(s,\cdot)$, with coupled states
$(X,Y)$. Its $i$th coordinate marginal is a coupling of
$p_{t,i}^\star(\cdot\mid s)$ and $p_{t,i}^q(\cdot\mid s)$. By the coupling characterization of
total variation (Lemma~\ref{lem:G:maxcoupling}),
$\Pr_\gamma[X_i\neq Y_i]\ge\delta_{t,i}(s)$. Therefore
\[
\Ex_\gamma[d(X,Y)]
=\sum_i\lambda_i\Pr_\gamma[X_i\neq Y_i]
\ge\sum_i\lambda_i\delta_{t,i}(s).
\]
Under the independent coordinatewise maximal coupling, equality holds for every coordinate:
$\Pr[X_i\neq Y_i]=\delta_{t,i}(s)$. Its expected weighted-Hamming cost is therefore the lower
bound above, proving optimality.
\end{proof}

\begin{remark}[Two one-step objectives]
\label{rem:G:cmcscope}
The product CMC is Wasserstein-optimal for the additive one-step cost $d$, but it is not generally
a maximal coupling of the \emph{joint} transition laws, which instead minimizes the $0$--$1$ event
$\Pr[X\neq Y]$. Nor is the iterated product construction claimed to be a Wasserstein-optimal
\emph{path} or terminal coupling. The direct and first-divergence routes can therefore prefer
different one-step couplings.
\end{remark}



\subsection{Proof of Proposition~\ref{prop:lifted} (lifted sample-then-select coupling)}
\label{app:proofs:lifted}

\noindent\textbf{Proposition~\ref{prop:lifted} (restated).}\par
\noindent \emph{For any coupling $\pi_t$ of the proposal laws, the pushforward pair
$(\Phi_t^\star(s,\xi^\star),\Phi_t^q(s,\xi^q))$ is a coupling of the actual transitions. Its
split probability is at most the sum of proposal mismatch probability and the probability that an
equal proposal is mapped to different next states. Under product proposal CMC the first term is
$1-\prod_i(1-\delta_{t,i}^{\mathrm{prop}})$. Iteration yields a valid path coupling, and when the
pushed-forward transitions are monotone copy-only the first-divergence and bounded-functional
results apply.}

\begin{proof}
Let $(\xi^\star,\xi^q)\sim\pi_t(\cdot,\cdot\mid s)$. For any $A\subseteq\Xs$,
\begin{align*}
\Pr\!\left[\Phi_t^\star(s,\xi^\star)\in A\right]
&=\nu_t^\star\!\left(\{\xi:\Phi_t^\star(s,\xi)\in A\}\mid s\right)
=P_t(s,A),\\
\Pr\!\left[\Phi_t^q(s,\xi^q)\in A\right]
&=\nu_t^q\!\left(\{\xi:\Phi_t^q(s,\xi)\in A\}\mid s\right)
=Q_t(s,A).
\end{align*}
Thus the pushforward pair has the required marginals. For the split event, partition according to
whether the proposals agree:
\[
\{\Phi_t^\star(s,\xi^\star)\ne\Phi_t^q(s,\xi^q)\}
\subseteq
\{\xi^\star\ne\xi^q\}
\cup
\{\xi^\star=\xi^q,\,\Phi_t^\star(s,\xi^\star)\ne\Phi_t^q(s,\xi^q)\}.
\]
Taking probabilities proves the channel upper bound. If the proposal laws are coordinate products
and $\pi_t$ is the independent product of coordinatewise maximal couplings, proposal equality has
probability $\prod_i(1-\delta_{t,i}^{\mathrm{prop}})$, giving the stated first term.

To form a path coupling, use the lifted one-step coupling whenever the two current states agree and
any fixed valid coupling otherwise. The same induction used in the proof of
Theorem~\ref{thm:cmc} shows that the two path marginals are the FP and quantized chains. If each
marginal transition is monotone copy-only, Theorem~\ref{thm:fs} and its corollaries apply to this
valid path coupling without any factorization assumption on the pushed-forward transition.
\end{proof}

\subsection{Proof of Corollary~\ref{cor:det-limit} (temperature-zero limit)}
\label{app:proofs:det-limit}

\paragraph{Corollary~\ref{cor:det-limit} (restated).}
\emph{If the two proposal laws are Dirac masses, the proposal coupling is unique; coupling-choice
comparisons disappear, the next-state pair is deterministic, and under monotone copy-only decoding
the deterministic action decomposition and margin lemmas apply.}

\begin{proof}
A coupling of $\delta_{\xi_t^\star}$ and $\delta_{\xi_t^q}$ must assign probability one to the
single pair $(\xi_t^\star,\xi_t^q)$, so there is no coupling freedom. Applying the deterministic maps
$\Phi_t^\star$ and $\Phi_t^q$ produces the deterministic next-state pair
$(\Phi_t^\star(s,\xi_t^\star),\Phi_t^q(s,\xi_t^q))$. If the resulting actions are monotone
copy-only, they fall exactly under the deterministic setting of Theorem~\ref{thm:exact}; the
margin-necessity lemmas then apply to token and policy boundary crossings as stated.
\end{proof}

\subsection{Proof of Theorem~\ref{thm:fixedsuite} (fixed-suite simultaneous certificate)}
\label{app:proofs:stat}

\paragraph{Theorem~\ref{thm:fixedsuite} (restated).}
\emph{Fix prompts $c_{1},\dots,c_{m}$, a predeclared finite family of cells
$j\in\mathcal{J}$ with $K=|\mathcal{J}^{\uparrow}|+|\mathcal{J}^{\downarrow}|$ one-sided
statements, and for each cell $r_{j}$ draws
$X_{a,s,j}\in[0,1]$, $a\le m$, $s\le r_{j}$, independent across $(a,s)$ within the cell
(arbitrary dependence across cells). Let
$\hat\mu_{j}=\frac{1}{mr_{j}}\sum_{a,s}X_{a,s,j}$ and
$\mu_{j}=\frac{1}{m}\sum_{a}\Ex[X_{a,1,j}]$. Then with probability at least $1-\delta$,
simultaneously for every upper cell
$\mu_{j}\le\hat\mu_{j}+\sqrt{\log(K/\delta)/(2mr_{j})}$ and for every lower cell
$\mu_{j}\ge\hat\mu_{j}-\sqrt{\log(K/\delta)/(2mr_{j})}$.}

\begin{proof}
Fix an upper cell $j$ and write $n=mr_{j}$. The variables
$\{X_{a,s,j}\}$ are independent, $[0,1]$-valued, and not necessarily identically distributed;
their means average to
$\frac{1}{n}\sum_{a,s}\Ex[X_{a,s,j}]=\frac{1}{m}\sum_{a}\Ex[X_{a,1,j}]=\mu_{j}$, using that
draws within a prompt are identically distributed. Hoeffding's inequality for independent bounded
variables gives, for every $\varepsilon>0$,
\[
\Pr\big[\hat\mu_{j}\le\mu_{j}-\varepsilon\big]\ \le\ \exp\!\big(-2n\varepsilon^{2}\big).
\]
Choosing $\varepsilon_{j}=\sqrt{\log(K/\delta)/(2mr_{j})}$ makes the right side $\delta/K$. The
symmetric tail handles each lower cell with the same $\varepsilon_{j}$. A union bound over the $K$
one-sided statements bounds the probability that \emph{any} of them fails by $\delta$; on the
complementary event all $K$ inequalities hold simultaneously. No assumption is made about
dependence \emph{across} cells: the union bound is valid under arbitrary joint laws.
\end{proof}

\subsection{Proof of Corollary~\ref{cor:fixedsuite} (fixed-suite CMC certificate)}
\label{app:proofs:fixedsuite-cor}

\paragraph{Corollary~\ref{cor:fixedsuite} (restated).}
\emph{If for every prompt $a$ and cell $j$,
$W_{a,j}\le\Ex[D_{a,1,j}]\le\Ex[B_{a,1,j}]$, then simultaneous applications of
Theorem~\ref{thm:fixedsuite} to the preregistered $D$ and $B$ cells yield
$\overline W_j\le\min\{1,\operatorname{UCB}(\widehat D_j),
\operatorname{UCB}(\widehat B_j)\}$ for all cells.}

\begin{proof}
Average the prompt-wise deterministic inequalities to obtain
$\overline W_j\le m^{-1}\sum_a\Ex[D_{a,1,j}]$ and
$\overline W_j\le m^{-1}\sum_a\Ex[B_{a,1,j}]$. On the simultaneous event of
Theorem~\ref{thm:fixedsuite}, these two averages are at most their corresponding UCBs for every
cell. The metric diameter supplies the additional bound $\overline W_j\le1$; taking the minimum
preserves validity. Since the event is simultaneous over the preregistered horizon grid, any
deterministic prefix-closed selection rule based on the realized UCBs inherits the same coverage.
\end{proof}

\begin{remark}[Post-hoc horizon selection is licensed by simultaneity]
\label{rem:G:posthoc}
The certified-horizon rule in the main text selects, after observing the data, the largest
prefix-closed horizon whose UCBs all fall below a tolerance. Because the guarantee of
Theorem~\ref{thm:fixedsuite} is simultaneous over the entire predeclared family, any selection rule that is a
deterministic function of the realized UCBs inherits the same $1-\delta$ coverage; no correction
for selection is needed beyond the union bound already paid. This is exactly why the family
$\mathcal{J}$ must be declared before data collection, as done in the protocol.
\end{remark}

\begin{remark}[Independence premise]
\label{rem:G:indep}
Applying Theorem~\ref{thm:fixedsuite} to the CMC draws treats the $r_{j}$ event-addressed pseudorandom streams per
prompt as independent draws. This is an implementation premise about the PRNG, stated as such in
the statistical-certification and limitations sections; the collision-freedom and reproducibility checks of the audit support but do
not prove it. The conservative prompt-block sensitivity, which averages within prompts first ($n=m$), does not rely on within-prompt independence; it still assumes independence across prompt blocks.
\end{remark}

\subsection{Proof of Theorem~\ref{thm:population} (prompt-population concentration)}
\label{app:proofs:population}

\paragraph{Theorem~\ref{thm:population} (restated).}
\emph{Let $C_1,\dots,C_m\stackrel{\mathrm{iid}}{\sim}\Psrc$. For each preregistered cell $j$,
let $X_{a,1,j},\dots,X_{a,r_j,j}\in[0,1]$ have common conditional mean $g_j(C_a)$, and suppose
the prompt blocks $(C_a,X_{a,1:r_j,j})$ are independent across $a$ within the cell. Define
$Y_{a,j}=r_j^{-1}\sum_sX_{a,s,j}$,
$\widehat\nu_j=m^{-1}\sum_aY_{a,j}$, and
$\nu_j=\Ex_{C\sim\Psrc}[g_j(C)]$. For a preregistered family of $K$ one-sided statements, with
probability at least $1-\delta$, simultaneously for every upper cell
$\nu_j\le\widehat\nu_j+\sqrt{\log(K/\delta)/(2m)}$, with the analogous lower bound for every
lower cell. Dependence across cells is unrestricted.}

\begin{proof}
Fix a cell $j$. Since each $X_{a,s,j}$ lies in $[0,1]$, the block average $Y_{a,j}$ also lies in
$[0,1]$. Conditional on $C_a$,
\[
\Ex[Y_{a,j}\mid C_a]
=\frac1{r_j}\sum_{s=1}^{r_j}\Ex[X_{a,s,j}\mid C_a]
=g_j(C_a),
\]
where no independence among the repeated draws inside one prompt block is needed. Taking expectation
over the IID prompt draw gives
$\Ex[Y_{a,j}]=\Ex_{C\sim\Psrc}[g_j(C)]=\nu_j$. By assumption the prompt blocks are independent
across $a$, so $Y_{1,j},\dots,Y_{m,j}$ are independent $[0,1]$-valued random variables with common
mean $\nu_j$. One-sided Hoeffding therefore gives
\[
\Pr\!\left[\widehat\nu_j\le\nu_j-\varepsilon\right]
\le \exp(-2m\varepsilon^2)
\]
for an upper cell, and the symmetric upper-tail inequality yields the lower-cell statement. Taking
$\varepsilon=\sqrt{\log(K/\delta)/(2m)}$ makes the failure probability at most $\delta/K$ for
one statement. A union bound over the $K$ preregistered one-sided statements proves simultaneous
coverage. Dependence across cells is irrelevant to the union bound.
\end{proof}

\subsection{Proof of Corollary~\ref{cor:population} (prompt-population CMC certificate)}
\label{app:proofs:population-cor}

\paragraph{Corollary~\ref{cor:population} (restated).}
\emph{Suppose for every prompt $c$ and cell $j$,
$W_j(c)\le\Ex_U[D_j(c,U)]\le\Ex_U[B_j(c,U)]$. Under the sampling conditions of
Theorem~\ref{thm:population}, simultaneous UCBs based on prompt-block averages satisfy
$\Ex_C[W_j(C)]\le\min\{1,\operatorname{UCB}(\widehat D_j),
\operatorname{UCB}(\widehat B_j)\}$ for every preregistered cell.}

\begin{proof}
The prompt-wise coupling inequalities imply, after averaging over $C\sim\Psrc$,
\[
\Ex_C[W_j(C)]
\le \Ex_C\Ex_U[D_j(C,U)]
\le \Ex_C\Ex_U[B_j(C,U)].
\]
Apply Theorem~\ref{thm:population} to the prompt-block averages of the repeated $D$ draws and,
separately, to those of the repeated $B$ draws. On the simultaneous event, the two population
expectations on the right are bounded by their corresponding UCBs in every cell. The normalized
Hamming diameter adds the deterministic bound $\Ex_C[W_j(C)]\le1$, so taking the minimum preserves
validity. Since the event is simultaneous on the preregistered horizon grid, any deterministic
prefix-closed selection rule based on these UCBs inherits the same coverage, by the same argument as
Remark~\ref{rem:G:posthoc}.
\end{proof}

\begin{remark}[Population and fixed-suite targets are different]
Theorem~\ref{thm:fixedsuite} conditions on a fixed collection of prompts and may use every
independent coupling draw as a bounded observation. Theorem~\ref{thm:population} instead treats the
IID prompt block as the sampling unit; repeated conditional draws are averaged within that block.
The resulting estimands and Hoeffding radii therefore must not be pooled or interchanged.
\end{remark}
